\cleardoubleoddpage
\chapter{Discussion}
\label{chap:discussion}

\section{Answer to the Research Question}

The cache-normalized Transformer is the best learned model on validation flux RMSE. In the
retrospective seed-52 test realization, its flux RMSE is 11.4255, compared with 11.9889 for
decoded latent persistence and 3.7208 for direct persistence of the observed flux. Applying the
frozen diagnostic head to true future latents gives 12.0712, so the learned model does not beat the
appropriate observed-diagnostic baseline and the comparison does not establish superiority.

The run establishes that the complete data, representation, cache, sequence, diagnostic, and
provenance path can be executed with a strict trajectory-level split. It also shows that the
observed diagnostic is a much stronger reference than decoded latent persistence, and that a
complex temporal architecture does not by itself demonstrate useful forecasting. The oracle
result identifies diagnostic decoding as a limiting factor under the present encoder and budget.

\section{Interpretation of the Test Difference}

The Transformer and observed-persistence flux RMSE values differ by 7.7047 in the trajectory-balanced
headline, with the learned error 3.07 times as large. The five trajectory-level differences
(Transformer minus observed persistence) have mean 7.7902 and a descriptive 95\% interval from
$-0.6813$ to $16.2616$. The current design has only one sequence-training seed and five
trajectories, so this interval is descriptive and does not establish repeatability.

The Transformer has lower latent error than decoded latent persistence, but that comparison shares
the diagnostic head and is secondary to the observed-diagnostic baseline. The modest spectral
differences occur without a meaningful change in aggregate shape correlation. These disagreements
justify the metric hierarchy: observed-flux RMSE is primary, while latent error, MAE, and spectra
diagnostics explain the failure mode.

The learned model also has lower latent MSE than decoded latent persistence on the retrospective
test. Because diagnostic heads are nonlinear and task-specific, latent and observed-diagnostic
rankings need not match. Cache normalization improved both learned architectures on validation,
but matched multi-seed runs are needed before treating that effect as repeatable.

\section{Importance of the Integrity Audit}

The most consequential finding was methodological rather than architectural. Different trajectory
split seeds had previously been used for encoder training and downstream evaluation. A trajectory
can then be called ``test'' downstream even though its snapshots were used to train the encoder,
which violates end-to-end held-out evaluation. The resulting metric may still measure sequence-model
behaviour on a fixed cache, but it cannot support a claim about generalization of the complete
surrogate.

Withdrawing the earlier result is necessary scientific hygiene. The corrected system now binds
every cache and checkpoint to its split seed, records exact trajectory manifests, fits normalization
statistics on training trajectories only, and rejects incompatible artifacts before evaluation.
These controls are more valuable than retaining a numerically stronger but invalid headline. They
also make future negative or positive results easier to audit.

\section{Comparison with Related Work}

GyroSwin predicts full five-dimensional dynamics and targets physically motivated field validation
at a larger architectural scale \citep{paischer2025gyroswin}. The present system instead evaluates a
128-dimensional diagnostic-oriented state without a decoder. The two approaches therefore answer
different questions: field fidelity versus compact diagnostic prediction.

A direct accuracy comparison remains invalid without the same data revision, trajectory split,
temporal sampling, target definitions, normalization, forecast horizon, and aggregation. The
internal result does not imply that latent surrogates in general are superior to full-field
models, nor does it provide evidence about GyroSwin performance. It only constrains the specific
encoder, learned sequence candidates, data subset, and budget evaluated here.

\section{Limitations}

The study has several limitations:

\begin{itemize}
  \item The medium cache contains 51 trajectories from one preprocessing pipeline. Five validation
        and five test trajectories provide limited independent coverage of the underlying regime.
  \item Each learned candidate is trained once, and the GRU and Transformer use different fixed
        training seeds. Multiple matched seeds are required to separate architecture from
        initialization variability and assess repeatability.
  \item The eight-step horizon establishes short-term behaviour only. It does not demonstrate
        long-term stability or error saturation.
  \item The pointwise encoder uses aggressive subsampling and global pooling. It may discard spatial
        relationships that a temporal model would need to beat persistence.
  \item Diagnostic heads evaluate selected scalar and spectral targets. They cannot establish that
        the latent state retains all physically relevant degrees of freedom.
  \item Relative spectral errors are large and sensitive to small target norms. Their aggregate
        values combine two spectra with different scales and must be interpreted per target.
  \item The corrected run retrains the seed-52 encoder, cache, and sequence checkpoints from
        recorded source commits. Raw data and large artifacts are not distributed in Git, so exact
        reproduction requires authorized dataset access.
  \item Runtime speedup over the original gyrokinetic solver is not measured end to end and is not
        claimed.
\end{itemize}

\section{Future Work}

The next protocol uses matched training seeds 52--56. If at least ten previously uninspected
trajectories are available, their manifest becomes the final test after architecture,
normalization, budget, and checkpoints are fixed on validation. Otherwise, the study uses five-fold
nested group cross-validation and calls the analysis retrospective
\citep{varma2006crossvalidation,cawley2010selectionbias}. The primary interval is a paired
hierarchical bootstrap over training seeds and trajectories rather than overlapping windows
\citep{saravanan2020hierarchical}.

Architecturally, the encoder deserves more attention than a wider sequence sweep. Structured
five-dimensional patch mixing or less aggressive subsampling may retain temporal cues that the
pointwise hierarchy loses. Diagnostic objectives should normalize spectra per target so that one
scale cannot dominate the encoder loss. Longer contexts and horizons are appropriate only after a
learned model reliably exceeds persistence across seeds in the short-horizon setting.

If a complete GyroSwin checkpoint and evaluation bundle becomes available, a matched benchmark can
be designed after reproducing its native protocol. A decoder could separately test full-field
information retention, but that would extend the research question rather than repair the present
diagnostic-surrogate experiment.
